\scanpage{75}\mspage{J}
car $>0$. Cela prouve $r=0$, et achève la
démonstration.

\uline{Corollaire} 1 Invariance birationnelle de $\underline{\mathrm{Pic}}^\tau_{X/k}$
pour $X$ propre et simple sur $k$.

\uline{Corollaire} 2 Soit $K$ une extension
\deleted{régulière}\inserted{séparable primaire} de type fini de $k$, admettant
un modèle \deleted{régulier \unsure{et}} propre et \uline{simple}
sur $k$. Alors $\varinjlim_X\underline{\mathrm{Pic}}^\tau_{X/k}$, où $X$ parcourt
les Modèles propres et normaux \inserted{``sur'' $k$} de $K$
existe au sens \unsure{trivial|strict} [le système \deleted{projectif}\inserted{inductif}
est essentiellement constant]
\margin{\uline{N.B.} Si $k$ n'est pas parfait, et $K$ est de dim $1$ sur $k$, \deleted{\unsure{sur}} c'est une \unsure{restriction|condition} un peu idiote !.}
% slanted

\uline{Question}. Le cor. 2 reste-t-il valable
en supposant seulement que $K$ admette
un modèle ``normal sur $k$'' ?? \deleted{Il y aurait} \deleted{\unsure{On}}\inserted{\unsure{Cela|Ceci|Car|le cas}}
\deleted{du moins} équivaudrait à la conjonction des
assertions suivantes

(i) $\dim\underline{\mathrm{Pic}}_{X/k}$ reste majoré

(ii) ordre $\mathrm{Tors}\,\underline{\mathrm{Pic}}(X/k)$ reste majoré

(iii) $\dim H^1(X,\underline{O}_X)$ reste majoré.
% Tors Pic and the H^1 expression in the last two lines are individually encircled.
\margin{N.B. (i) résulte de (iii)}
% archive sheet number: 74
